% ===== PRÉAMBULE CANONIQUE — à coller VERBATIM en tête de CHAQUE .tex =====
\documentclass[11pt,a4paper]{article}
\usepackage[utf8]{inputenc}\usepackage[T1]{fontenc}\usepackage{lmodern}
\usepackage[french]{babel}
\usepackage{amsmath,amssymb,mathtools}\usepackage{icomma}
\usepackage[version=4]{mhchem}          % \ce{...} : équations & formules
\usepackage{siunitx}\sisetup{locale=FR} % \qty{}{}, \si{}, \ang{}
\usepackage{chemfig}                    % \chemfig{...} : structures développées
\usepackage{xcolor}\usepackage{colortbl}
\usepackage{tikz}\usepackage{pgfplots}\pgfplotsset{compat=1.18}
\usepackage[most]{tcolorbox}\usepackage{enumitem}\usepackage{array,booktabs}
\usepackage[a4paper,margin=2.2cm]{geometry}
\usetikzlibrary{arrows.meta,calc,angles,quotes,positioning,patterns,backgrounds,shapes.geometric}
\definecolor{primary}{RGB}{222,49,99}\definecolor{primarydark}{RGB}{150,22,55}
\definecolor{defcol}{RGB}{0,114,178}\definecolor{methcol}{RGB}{52,92,148}
\definecolor{excol}{RGB}{0,135,90}\definecolor{attncol}{RGB}{200,40,55}
\tcbset{boxbase/.style={enhanced,breakable,boxrule=0.9pt,arc=2.5pt,
  left=3mm,right=3mm,top=2.3mm,bottom=2.3mm,fonttitle=\bfseries,coltitle=white}}
% --- boîtes TUTORIEL (noms EXACTS, ne pas renommer) ---
\newtcolorbox{cle}[1][]{boxbase,colback=defcol!6,colframe=defcol,
  title={\textsc{À retenir}\ifx\relax#1\relax\relax\else\textmd{ : \itshape #1}\fi}}
\newtcolorbox{methode}[1][]{boxbase,colback=methcol!7,colframe=methcol,sharp corners,
  title={\textsc{Méthode}\ifx\relax#1\relax\relax\else\textmd{ --- \itshape #1}\fi}}
\newtcolorbox{exemplebox}[1][]{boxbase,colback=excol!5,colframe=excol,
  title={\textsc{Exemple}\ifx\relax#1\relax\relax\else\textmd{ : \itshape #1}\fi}}
\newtcolorbox{piege}[1][]{boxbase,colback=attncol!6,colframe=attncol,
  title={\textsc{Piège}\ifx\relax#1\relax\relax\else\textmd{ : \itshape #1}\fi}}
\newtcolorbox{remarque}[1][]{boxbase,colback=black!4,colframe=black!45,
  title={\textsc{Remarque}\ifx\relax#1\relax\relax\else\textmd{ : \itshape #1}\fi}}
% --- boîtes EXERCICE / CORRIGÉ (noms EXACTS, auto-comptées) ---
\newtcolorbox[auto counter]{exo}[1][]{boxbase,colback=primary!4,colframe=primary,
  title={\textsc{Exercice}~\thetcbcounter\ifx\relax#1\relax\relax\else\textmd{ --- \itshape #1}\fi}}
\newtcolorbox[auto counter]{corrige}[1][]{boxbase,colback=excol!4,colframe=excol,
  title={\textsc{Corrigé}~\thetcbcounter\ifx\relax#1\relax\relax\else\textmd{ --- \itshape #1}\fi}}
\newcommand{\R}{\mathbb{R}}\newcommand{\N}{\mathbb{N}}\newcommand{\Z}{\mathbb{Z}}\newcommand{\ds}{\displaystyle}
% (un \usepackage{fancyhdr}+en-tête est possible ; il est ignoré côté web)
% =========================================================================
\begin{document}
\sisetup{per-mode=symbol}
\begin{corrige}[du nombre de CS à la précision relative]
On lit directement les ordres de grandeur.
\begin{enumerate}[label=\alph*)]
  \item $2$ CS $\Rightarrow$ incertitude relative $\sim 1\%$.
  \item $3$ CS $\Rightarrow \sim 0,1\%$.
  \item $4$ CS $\Rightarrow \sim 0,01\%$.
\end{enumerate}
\end{corrige}

\begin{corrige}[quel instrument est le plus précis ?]
\begin{enumerate}[label=\alph*)]
  \item Éprouvette : $\dfrac{1}{25}=0,04=\boxed{4\,\%}$. \quad Burette : $\dfrac{0,05}{24,50}=0,00204\approx\boxed{0,20\,\%}$.
  \item La \textbf{burette} est bien plus précise ($0,20\%\ll 4\%$). C'est cohérent avec le nombre de CS :
        $25$ n'a que $2$ CS, tandis que $24,50$ en a $4$.
\end{enumerate}
\end{corrige}

\begin{corrige}[une balance déréglée]
Moyenne des cinq mesures : $\dfrac{50,01+49,99+50,00+50,02+50,00}{5}=\qty{50,00}{\gram}$.
\begin{enumerate}[label=\alph*)]
  \item \textbf{Oui}, elle est fidèle : les résultats sont très serrés (écart maximal $\qty{0,03}{\gram}$).
  \item \textbf{Non}, elle n'est pas juste : la moyenne $\qty{50,00}{\gram}$ diffère de la vraie valeur
        $\qty{48,00}{\gram}$ (biais systématique de $+\qty{2,00}{\gram}$).
  \item Il faut \textbf{étalonner} la balance (régler le zéro, ou retrancher le biais de $\qty{2,00}{\gram}$).
        Ajouter des CS n'y changerait rien : les CS ne corrigent pas un biais.
\end{enumerate}
\end{corrige}

\begin{corrige}[une écriture incohérente]
\begin{enumerate}[label=\alph*)]
  \item La valeur $5,2$ n'est écrite qu'au \emph{dixième} de gramme, mais l'incertitude $\qty{0,01}{\gram}$
        prétend une fiabilité au \emph{centième}. Valeur et incertitude doivent porter sur le même rang.
  \item Deux écritures cohérentes : $(5,20\pm0,01)\ \si{\gram}$ (balance au centième, $3$ CS) \emph{ou}
        $(5,2\pm0,1)\ \si{\gram}$ (balance au dixième, $2$ CS).
\end{enumerate}
\end{corrige}

\begin{corrige}[de la précision relative au nombre de CS]
On lit la table à l'envers.
\begin{enumerate}[label=\alph*)]
  \item $\sim 1\% \Rightarrow 2$ CS.
  \item $\sim 0,01\% \Rightarrow 4$ CS.
\end{enumerate}
\end{corrige}

\end{document}
